\title[Finite failure schemes and quadratic pencils]{Finite failure schemes and the reconstruction of quadratic pencils}
\author{Qian Qi}
\date{September 24, 2026}
\subjclass[2020]{14M12, 14B05, 14D20, 13C40, 62R01}
\keywords{Intrinsic reconstruction, quadratic pencils, Fitting schemes, finite neighbourhoods, spectral sheaves, likelihood correspondence}
\hypersetup{pdftitle={Finite failure schemes and the reconstruction of quadratic pencils, revision 143},pdfauthor={Qian Qi}}
\begin{abstract}
Every quadratic pencil in dimension \(n\ge3\) is determined up to
projective equivalence by the abstract \((n^2+2n-4)\)-th neighbourhood
of the socle Grassmannian in its multiplication-failure scheme.
Every lower-order neighbourhood is independent of the pencil.
We prove this by recovering the oriented tensor cone and its full
Pluecker coefficient line intrinsically.  The first relation spaces
embed the entire pencil parameter scheme and give universal finite
neighbourhoods compatible with arbitrary complex base change.
For regular pencils they recover spectral sheaves, their Fitting
incidence schemes, and the reciprocal critical correspondence.
At fixed discriminant and reduced rank data, the finite invariant
realizes and distinguishes every spectral dominance transition.
For definite real pencils, spectral projectors give a congruence-covariant
open family of data on which the critical correspondence is finite
etale and all its points are real, simple, and globally labelled.
The full web, contraction, polarized, and boundary theory is retained
in the accompanying supplement.
\end{abstract}
\maketitle
